\documentclass[12pt]{article}
\usepackage[utf8]{inputenc}
\usepackage[T1]{fontenc}
\usepackage{hyperref}
\title{Kathu Language Resources for K-12 Curriculum Development}
\author{Compiled by Claude Sonnet 5.5 (claude-sonnet-5-5) for Andr\'as Kornai}
\date{}
\begin{document}
\maketitle

\section{Introduction}
Kathu (ISO 639-3: ykt; Glottocode kath1251) is a Sino-Tibetan language of the Lolo-Burmese branch (Mondzish). It is spoken in Balong, Nanping Township, Guangnan County, Yunnan, China, and is also called Gasu or Kasu. The speakers are locally called White Yi. Endangered Languages reports about 5,000 speakers and rates the language as threatened. A 2004 estimate puts the total above 7,000 across Guangnan and Jinping counties. Joshua Project classifies the Kathu as Gelao and the language as Tai-Kadai. This conflicts with the other sources, so it should be treated with caution and checked. The search found no evidence that Kathu is taught in schools.

\section{Primary Educational Resources}
The search found no primers, textbooks or school materials in Kathu.

\section{Linguistic and Reference Materials}
\begin{itemize}
\item Wikipedia article on the Kathu language (classification, location, speaker numbers): \url{https://en.wikipedia.org/wiki/Kathu_language}
\item Austronesian Basic Vocabulary Database entry for Kathu, which lists a basic word list: \url{https://abvd.eva.mpg.de/austronesian/language.php?id=965}
\item Verbix page on Kathu as an endangered language: \url{https://docs.verbix.com/EndangeredLanguages/Kathu}
\end{itemize}

\section{Religious and Community Resources}
\begin{itemize}
\item Joshua Project profile of the Kathu in China. Its ethnolinguistic classification disagrees with the other sources: \url{https://joshuaproject.net/people_groups/20602/CH}
\end{itemize}
The search found no Kathu scripture translation and no community organisation.

\section{Digital and Audio Resources}
\begin{itemize}
\item Endangered Languages Project entry for Kathu: \url{https://endangeredlanguages.com/lang/5193}
\end{itemize}
The search found no audio or video recordings.

\section{Curriculum Development Context}
Very little material exists in the search results, and none of it is pedagogical. A curriculum would need new documentation work, including a practical orthography, word lists and recordings. It should be developed with the speaker community. The classification dispute and the differing speaker counts should be resolved through fieldwork or the linguistic literature.

\end{document}